%%==========================================================================%%
%% validation.tex -- "Technical Validation". Quality-scoring template modeled %%
%% on the reference paper, paraphrased/genericized. REWRITE before submission.%%
%% [TBD: ...] = fill in.                                                       %%
%%==========================================================================%%

\section*{Technical Validation}\label{sec:validation}

\subsection*{Expert quality assessment}\label{subsec:imgqc}
% Reusable: expert (radiologist/radiographer) scan-by-scan quality rating on an
% ordinal Likert scale (1-5). Adapt scale wording to whole-body as needed.
Each examination underwent a scan-by-scan quality assessment by [TBD: an expert
reviewer / board-certified radiologist], who rated diagnostic image quality on a
5-point Likert scale ranging from 1 (non-diagnostic) to 5 (perfect quality;
Table~\ref{tab:qc}). [TBD: report the distribution of scores and how qc\_status
was assigned, and note any examinations excluded as a result.]

%%-- Quality-scoring table placeholder (cf. Table 2 of ref) ----------------%%
\begin{table}[h]
\caption{Image-quality scoring scale used by the expert reviewer
(placeholder --- adapt wording for whole-body MRI).}\label{tab:qc}
\begin{tabular}{@{}cl@{}}
\toprule
Score & Description \\
\midrule
5 & Perfect scan, no artefacts. \\
4 & Mild artefacts, fully diagnostic. \\
3 & Mild--moderate artefacts; repeat if an abnormality were noticed. \\
2 & Substantial artefacts; large abnormalities discernible but repeat needed. \\
1 & Completely non-diagnostic. \\
\botrule
\end{tabular}
\end{table}

\subsection*{Inter-reader agreement}\label{subsec:agreement}
% Because labels are keyed by reader, the multiply-read subset supports
% inter-reader agreement: Dice/IoU for segmentations, Cohen's/Fleiss' kappa for
% categorical labels. Strong validation content.
For the subset of examinations annotated by multiple readers, we report
inter-reader agreement to quantify label reliability --- [TBD: Dice/IoU for
segmentations and Cohen's/Fleiss' $\kappa$ for categorical labels] --- broken
down by anatomy where appropriate.

\subsection*{Metadata and data integrity}\label{subsec:integrity}
% Reusable concept: confirm metadata completeness/consistency + file integrity.
We verified the integrity and internal consistency of the release: [TBD: schema
validation of the metadata tables, referential consistency with the relational
index, and per-file checksums (e.g., SHA-256)]. [TBD: note the repository's
preflight validation if applicable --- e.g., the PhysioNet validation tool.]

\subsection*{Radiology report de-identification}\label{subsec:reportqc}
% Validation that scrubbed reports are PHI-free.
Where radiology reports are released, we audited the de-identification to confirm
removal of protected health information using [TBD: manual review of a sample
and/or automated PHI detection]; [TBD: report the outcome].

\subsection*{Baseline analysis}\label{subsec:baseline}
% Optional: baseline model / sanity check demonstrating fitness for purpose.
[TBD: optional baseline experiment demonstrating that the data are fit for the
intended whole-body anatomical classification/segmentation tasks.]
